\section*{Capítulo \ref{ch:b3:complex-spectra-magnetism} --- Espectros eletrônicos e magnetismo dos complexos}

\begin{solution}{exo:b3:complex-spectra-magnetism:1}
D: E$_g$ + T$_{2g}$ ($2 + 3 = 5$); F: A$_{2g}$ + T$_{1g}$ + T$_{2g}$ ($1 + 3 + 3 = 7$); G: A$_{1g}$ + E$_g$ + T$_{1g}$ + T$_{2g}$
($1 + 2 + 3 + 3 = 9$).
\end{solution}

\begin{solution}{exo:b3:complex-spectra-magnetism:2}
$d^1$ $^2$D $\to$ $^2$T$_{2g}$; $d^2$ $^3$F $\to$ $^3$T$_{1g}$; $d^3$ $^4$F $\to$ $^4$A$_{2g}$; $d^4$ $^5$D $\to$ $^5$E$_g$; $d^5$ $^6$S $\to$ $^6$A$_{1g}$; $d^6$ $^5$D
$\to$ $^5$T$_{2g}$; $d^7$ $^4$F $\to$ $^4$T$_{1g}$; $d^8$ $^3$F $\to$ $^3$A$_{2g}$; $d^9$ $^2$D $\to$ $^2$E$_g$.
\end{solution}

\begin{solution}{exo:b3:complex-spectra-magnetism:3}
\ce{Mn^{2+}} de spin alto ($n = 5$) $5.92\,\mu_B$, \ce{Fe^{2+}} (4) $4.90\,\mu_B$, \ce{Co^{2+}} (3) $3.87\,\mu_B$; \ce{Fe^{2+}} e
\ce{Co^{3+}} de spin baixo ($t_{2g}^6$): 0, diamagnéticos.
\end{solution}

\begin{solution}{exo:b3:complex-spectra-magnetism:4}
$\ce{[Mn(H2O)6]^{2+}}$ (proibida por spin e por Laporte, quase incolor) $<$ $\ce{[Co(H2O)6]^{2+}}$
(proibida por Laporte, vibrônica) $<$ $\ce{[CoCl4]^{2-}}$ (sem centro de simetria) $<$ \ce{MnO4-} (transferência
de carga permitida).
\end{solution}

\begin{solution}{exo:b3:complex-spectra-magnetism:5}
8500: $^3$A$_{2g}\to{}^3$T$_{2g}$, logo $\Delta_{\mathrm o} = \qty{8500}{cm^{-1}}$; 14100: $^3$T$_{1g}$(F); 24900: $^3$T$_{1g}$(P).
$15B = 39\,000 - 3 \times 8500$, $B = \qty{900}{cm^{-1}}$, $\beta = 900/1056 = 0.85$.
\end{solution}

\begin{solution}{exo:b3:complex-spectra-magnetism:6}
$\Delta_{\mathrm o} = \qty{17400}{cm^{-1}}$; resolvendo $7.5B + 1.5\Delta_{\mathrm o} - \frac12\sqrt{225B^2 - 18B\Delta_{\mathrm o} + \Delta_{\mathrm o}^2} = 24\,600$: $B = \qty{729}{cm^{-1}}$,
$\beta = 0.79$. $\nu_3 = \qty{38500}{cm^{-1}}$ (\qty{260}{nm}), no ultravioleta.
\end{solution}

\begin{solution}{exo:b3:complex-spectra-magnetism:7}
Os elétrons $d$ estão mais deslocalizados sobre os íons óxido das gemas do que sobre
as moléculas de água: ali, as ligações Cr--O são mais covalentes.
\end{solution}

\begin{solution}{exo:b3:complex-spectra-magnetism:8}
Forte ($e_g$ preenchido de modo desigual): $d^9$, $d^4$ de spin alto, $d^7$ de spin baixo. Nula: $d^3$ ($t_{2g}^3$, termo A) e
$d^6$ de spin baixo ($t_{2g}^6$). Fraca: $d^6$ de spin alto ($t_{2g}^4e_g^2$, termo T).
\end{solution}

\begin{solution}{exo:b3:complex-spectra-magnetism:9}
$d^7$ octaédrico de spin alto: termo fundamental $^4$T$_{1g}$, um termo T com \omterm{def:b3:complex-spectra-magnetism:moment}{contribuição orbital},
momento bem acima do valor de spin puro de $3.87\,\mu_B$. $d^7$ tetraédrico ($\equiv$ $d^3$ octaédrico): termo
fundamental $^4$A$_2$, momento orbital suprimido, próximo do valor de spin puro (um pouco acima, por
mistura com termos T excitados).
\end{solution}

\begin{solution}{exo:b3:complex-spectra-magnetism:10}
$\mu_{\mathrm{eff}} = 2.828\sqrt{1.87} = 3.87\,\mu_B$: três elétrons desemparelhados ($S = \frac32$).
\end{solution}

\begin{solution}{exo:b3:complex-spectra-magnetism:11}
$\Delta\chi_v = 3 \times 25.0/\num{400e6} = \num{1.88e-7}$; $\chi_{\mathrm m} = \num{1.88e-7}/\qty{10.0}{mol.m^{-3}} = \qty{1.88e-8}{m^3.mol^{-1}}$, isto é,
\qty{1.49e-3}{cm^3.mol^{-1}}; $\chi_{\mathrm m}T = 0.445$, $\mu_{\mathrm{eff}} = 1.89\,\mu_B$: um elétron desemparelhado.
\end{solution}

\begin{solution}{exo:b3:complex-spectra-magnetism:12}
$T_{1/2} = 15\,000/75 = \qty{200}{K}$. A \qty{250}{K}: $\Delta G = 15\,000 - 250 \times 75 = \qty{-3750}{J/mol}$, $K = \eu^{1.80} = 6.07$,
$x_{\mathrm{HS}} = 0.86$; $\chi_{\mathrm m}T = 0.86 \times 3.00 = \qty{2.6}{cm^3.K.mol^{-1}}$.
\end{solution}

\begin{solution}{pb:b3:complex-spectra-magnetism:1}
\textbf{1.} [Ar]$3d^3$.
\textbf{2.} $^4$F.
\textbf{3.} $^4$A$_{2g}$ + $^4$T$_{2g}$ + $^4$T$_{1g}$.
\textbf{4.} $^4$A$_{2g}$, o termo de $t_{2g}^3$, com os três elétrons nos orbitais mais baixos.
\textbf{5.} $^4$T$_{1g}$(P).
\textbf{6.} $^4$A$_{2g}\to{}^4$T$_{2g}$, $\to{}^4$T$_{1g}$(F), $\to{}^4$T$_{1g}$(P).
\textbf{7.} 561 e \qty{414}{nm}.
\textbf{8.} 597 e \qty{434}{nm}.
\textbf{9.} O rubi absorve o amarelo-esverdeado e o violeta e transmite o vermelho (com um pouco
de azul); as bandas da esmeralda, deslocadas para o vermelho, absorvem também o vermelho-alaranjado, e o verde
é transmitido.
\textbf{10.} Rubi \qty{17820}{cm^{-1}}, esmeralda \qty{16740}{cm^{-1}}.
\textbf{11.} $^4$T$_{2g}$ ($t_{2g}^2e_g$) e $^4$A$_{2g}$ ($t_{2g}^3$) têm a mesma energia de repulsão eletrônica; sua
diferença é a energia de promoção monoeletrônica $\Delta_{\mathrm o}$.
\textbf{12.} $B = (14\,305 - 547)/15 = \qty{917}{cm^{-1}}$.
\textbf{13.} \qty{614}{cm^{-1}}.
\textbf{14.} \qty{618}{cm^{-1}}.
\textbf{15.} 0.67 para ambas.
\textbf{16.} 29.0 e 27.1.
\textbf{17.} $\nu_3 = \qty{38500}{cm^{-1}}$ (\qty{260}{nm}) para o rubi, no ultravioleta, onde a absorção de
transferência de carga a encobre.
\textbf{18.} O sítio octaédrico do berilo é maior: ligações Cr--O mais longas dão um
desdobramento menor, que varia fortemente com a distância.
\textbf{19.} $10^7/695 = \qty{14390}{cm^{-1}}$, $23.4B$.
\textbf{20.} $^2$E$_g$ depende sobretudo de $C$: o nível medido mais alto significa que $C/B$ é maior,
cerca de 5.3 pela mesma diagonalização.
\textbf{21.} $^2$E$_g$ e $^4$A$_{2g}$ pertencem ambos a $t_{2g}^3$ e diferem apenas pela inversão de um spin: a ligação,
logo a geometria, é a mesma nos dois, e a transição não tem progressão
vibracional (uma raia 0--0).
\textbf{22.} Ela é proibida por spin.
\textbf{23.} Sua energia depende de $B$ e de $C$, quase nada de $\Delta_{\mathrm o}$ (uma linha plana no diagrama),
e $B$ é quase o mesmo nas duas gemas.
\textbf{24.} Os íons bombeados para as bandas largas relaxam e se acumulam no nível
$^2$E$_g$ de vida longa, de modo que pode haver mais íons nele do que no nível fundamental: uma
inversão de \omterm{def:b3:partition-functions:boltzmann}{população}.
\textbf{25.} \textbf{$\Delta_{\mathrm o}(\text{rubi}) - \Delta_{\mathrm o}(\text{esmeralda}) = \qty{1080}{cm^{-1}}$}: toda a diferença entre o vermelho e o
verde.
\end{solution}
